% ==========================================================================
% Planning reference: what each thesis chapter and section should contain.
% Heading levels match the actual thesis for easy cross-referencing.
% ==========================================================================

%--- ~3--5 pages ---%
\chapter{Introduction}\label{ch:introduction}
\markboth{Introduction}{}

\section{Motivation and Context}
\begin{itemize}
    \item Introduce the industrial relevance of autonomous robotic manipulation for textile recycling and waste sorting on conveyors
    \item Highlight the challenges of using tendon-driven robots: compliant design, complex actuation, need for simulation-based development
    \item Emphasize the role of simulation (NVIDIA Isaac Sim) and reinforcement learning as an enabling approach for developing control policies without physical hardware iteration
    \item Frame the project thesis as the foundational simulation and validation stage of a larger multi-agent textile sorting project
\end{itemize}

\section{Problem Statement and Objectives}
\begin{itemize}
    \item State the thesis problem: constructing, validating, and controlling a tendon-driven robot in simulation for conveyor-based pick-and-place tasks
    \item List the research objectives:
    \begin{enumerate}
        \item High-fidelity simulation environment (conveyor, drum, objects) in Isaac Sim 5.1.0
        \item Robot model integration and validation (5-DOF tendon-driven arm + base + gripper)
        \item RL-based reach and cube place tasks with PD and tendon-driven variants
        \item Tendon wrist transfer experiment on a standard industrial robot arm
    \end{enumerate}
    \item Scope delimitation: this thesis does NOT address cloth simulation, sorting (selectivity), multi-agent coordination, or sim-to-real transfer; these are explicitly deferred to the subsequent master thesis
\end{itemize}

\section{Thesis Outline}
\begin{itemize}
    \item Brief summary of chapters 2--7 for reader orientation
\end{itemize}

%--- ~15--20 pages ---%
\chapter{Background}\label{ch:background}
\markboth{Background}{}

\section{Reinforcement Learning for Robotic Manipulation}
\begin{itemize}
    \item MDP formalism: states, actions, transitions, rewards, discount factor
    \item Policy gradient methods relevant to continuous control: PPO (primary), SAC, DDPG (for comparison context)
    \item On-policy (PPO) vs.\ off-policy approaches: stability, sample efficiency, suitability for GPU-parallel simulation
    \item Policy network architectures: MLP actor-critic, observation normalization
    \item The skrl library as the RL framework used in this thesis
    \item Prior successes of RL in simulation for manipulation (OpenAI Dexterous Hand, QT-Opt, DexMV)
\end{itemize}

\section{Physics Simulation in Isaac Sim 5.1}
\begin{itemize}
    \item NVIDIA PhysX 5 as core solver, GPU-accelerated simulation
    \item Rigid body dynamics: TGS solver, articulated bodies, joints, contacts, friction
    \item Simulation parameters: time step, solver iterations, contact offsets
    \item Collision detection: mesh vs.\ convex decomposition vs.\ primitive shapes
    \item Brief mention of deformable simulation capabilities (PBD particle cloth, VBD) as context for future work --- not used in this thesis
\end{itemize}

\section{NVIDIA Isaac Sim and IsaacLab Architecture}
\begin{itemize}
    \item Isaac Sim 5.1.0: Omniverse platform, PhysX 5, RTX rendering, ROS 2 support
    \item IsaacLab: unified robotics learning framework, manager-based environment structure
    \item MDP workflow: observation managers, action managers, reward managers, event managers
    \item Comparison with alternative platforms (MuJoCo, PyBullet) to justify choice
\end{itemize}

\section{Tendon-Driven Robot Kinematics and Modeling}
\begin{itemize}
    \item Tendon-driven actuation: cables generate joint torques, advantages for compliant manipulation
    \item The specific robot: 2-DOF linear base + 3-DOF tendon arm (Klein 2023) + Robotiq 2F-140 gripper
    \item Jacobian-transpose torque mapping: 5 cable tensions to 3 joint torques
    \item Challenges of simulating tendon-driven systems in Isaac Sim; decision to use PD drives instead of PhysX tendon features
    \item Workspace analysis fundamentals: FK-based Monte Carlo sampling, Yoshikawa manipulability index, inverse condition number, reachability density
\end{itemize}

\section{Policy Learning Pipeline}
\begin{itemize}
    \item Observation space design for robot control tasks
    \item Action space: joint position targets, tendon tension targets
    \item Reward function design: dense shaping rewards, curriculum learning
    \item Episode termination and reset strategies
    \item GPU-parallel vectorized environments for training throughput
\end{itemize}

\section{Prior Work}
\begin{itemize}
    \item TODO: Related work on tendon-driven robot simulation and RL-based pick-and-place
\end{itemize}

\section{Robot and Scene Description Formats}
\begin{itemize}
    \item USD: extensible 3D scene graph for Isaac Sim
    \item URDF: robot kinematic/dynamic description
    \item URDF-to-USD import pipeline and modifications for the tendon arm
    \item Isaac Sim Robot Assembler for modular multi-component robots
\end{itemize}

%--- ~3 pages ---%
\chapter{Research Gap}\label{ch:research_gap}
\markboth{Research Gap}{}

\section{Critical Analysis of State of the Art}
\begin{itemize}
    \item Tendon-driven robots are studied in control theory but rarely integrated with modern RL simulation frameworks
    \item Most RL manipulation work uses standard rigid-joint robots (UR, Franka); tendon-driven arms are underrepresented
    \item No existing work combines a tendon-driven arm with IsaacLab's GPU-parallel RL training pipeline
    \item The transferability of tendon wrist mechanisms to standard arms is unexplored
\end{itemize}

\section{Research Gap and Justification}
\begin{itemize}
    \item Gap: validated simulation of a tendon-driven robot in Isaac Sim with RL-based control for manipulation tasks does not exist in the literature
    \item Scientific contribution: methodology for integrating tendon-driven actuation in a modern GPU-parallel RL framework; comparative validation against Klein (2023)
    \item Practical contribution: a modular, reusable simulation environment for tendon-robot research; wrist transferability results
\end{itemize}

\section{Refined Problem Definition}
\begin{itemize}
    \item Research question: How can a custom tendon-driven manipulator be faithfully modeled, validated, and controlled via RL in Isaac Sim for conveyor-based pick-and-place tasks, and does the tendon wrist mechanism transfer beneficially to a standard robot arm?
    \item Scope: simulation only; no physical hardware; no cloth; no sorting selectivity
    \item Expected contributions: validated robot model, RL pipeline, wrist transfer comparison
\end{itemize}

%--- ~15--18 pages ---%
\chapter{Methodology}\label{ch:methodology}
\markboth{Methodology}{}

\section{Overview of Approach}
\begin{itemize}
    \item High-level system diagram: Isaac Sim scene $\rightarrow$ IsaacLab environment $\rightarrow$ skrl PPO $\rightarrow$ learned policy
    \item Development stages: (1) environment and robot setup, (2) validation, (3) RL tasks, (4) wrist transfer experiment
\end{itemize}

\section{Simulation Environment Setup}
\begin{itemize}
    \item USD scene: ground plane, conveyor belt, plastic drum, dome lighting
    \item Object placement and randomization (initial positions, colours for curriculum)
    \item Physics parameters: time step, solver iterations
    \item Parallel environment configuration for GPU training
    \item Figure: annotated simulation scene
\end{itemize}

\section{Robot Model Integration}
\begin{itemize}
    \item 2-DOF articulated base: USD construction, prismatic joints, travel ranges
    \item 3-DOF tendon-driven arm: URDF import from Klein (2023), modifications, revolute joints with PD drives
    \item Tendon actuation modeling decision: why PhysX tendon features were not used; PD drive approximation
    \item Robotiq 2F-140 gripper: asset library integration, specifications
    \item Modular assembly via Robot Assembler: two-step procedure (base+arm, then arm+gripper)
    \item Collision geometry: convex hull generation, self-collision verification
    \item Joint limits table (Table~X)
    \item Figure: exploded view of assembled robot
\end{itemize}

\section{Control Gain Tuning and Validation}
\begin{itemize}
    \item Isaac Sim Gain Tuner: methodology for automated stiffness/damping tuning
    \item Step-response testing: sinusoidal and step commands, metrics (settling time, overshoot, steady-state error)
    \item Comparison with Klein (2023) PID controller parameters
    \item Gravity compensation for ceiling-mounted configuration
    \item Figure: Gain Tuner output (commanded vs.\ actual trajectories)
\end{itemize}

\section{Workspace Analysis}
\begin{itemize}
    \item Monte Carlo FK sampling: 200,000 samples, 4096 parallel environments
    \item Data collected: end-effector position, Jacobian, Yoshikawa manipulability, inverse condition number
    \item Voxelization and density aggregation
    \item Comparison with Klein (2023) task geometry
    \item Visualization: 3D scatter plots, 2D heatmap cross-sections
\end{itemize}

\section{Reinforcement Learning Setup}
\begin{itemize}
    \item \textbf{Reach Task:} observation space (joint states, target pose), action space (joint position targets), reward (pose distance + orientation), PD-driven and tendon-driven variants
    \item \textbf{Cube Place Task:} observation space (joint states, cube pose, drum pose, grasp status), action space, reward (approach + grasp + lift + place), curriculum (green-only $\rightarrow$ green+red)
    \item PPO hyperparameters (learning rate, discount, network architecture, number of environments)
    \item Reproducibility table: software versions, hardware
\end{itemize}

\section{Tendon Wrist Transfer Experiment}
\begin{itemize}
    \item Experimental design: UR10 with native rigid wrist vs.\ UR10 with tendon-driven wrist replacement
    \item Control variable: wrist type; all other parameters (RL config, task, scene) identical
    \item Tasks: reach accuracy, place success rate
    \item Metrics: final pose error, success rate, learning speed
    \item Integration procedure: wrist USD asset swapped via Robot Assembler
\end{itemize}

%--- ~12--15 pages ---%
\chapter{Results}\label{ch:results}
\markboth{Results}{}

\section{Robot Simulation Validation Results}
\begin{itemize}
    \item Step-response plots per joint: commanded vs.\ actual position
    \item Quantitative metrics: final stiffness/damping, settling time, overshoot, steady-state error
    \item Comparison with Klein (2023) PID results; discussion of differences
    \item Collision geometry validation across full joint range
\end{itemize}

\section{Workspace Analysis Results}
\begin{itemize}
    \item 3D reachable workspace envelope and volume estimate
    \item 2D cross-section heatmaps by reachability density and manipulability
    \item Yoshikawa manipulability and inverse condition number distributions
    \item Coverage of desired task workspace; practical implications for object placement
    \item Comparison with Klein (2023) workspace (without 2-DOF base)
\end{itemize}

\section{Reach Task Results}
\begin{itemize}
    \item Learning curves (reward vs.\ training steps) for PD and tendon-driven variants
    \item Final pose accuracy: position error, orientation error
    \item PD vs.\ tendon comparison: learning speed, final accuracy
    \item Computational cost: wall-clock time, GPU memory, environment steps
\end{itemize}

\section{Cube Place Task Results}
\begin{itemize}
    \item Learning curves with curriculum progression
    \item Success rate: cube in drum, colour selectivity (green placed, red ignored)
    \item PD vs.\ tendon comparison
    \item Qualitative policy behavior description with keyframe figures
\end{itemize}

\section{Tendon Wrist Transfer Results}
\begin{itemize}
    \item UR10 native wrist vs.\ UR10 + tendon wrist: reach accuracy comparison
    \item Place success rate comparison
    \item Learning curves comparison
    \item Statistical significance of observed differences
\end{itemize}

%--- ~5 pages ---%
\chapter{Discussion}\label{ch:discussion}
\markboth{Discussion}{}

\section{Evaluation of Outcomes}
\begin{itemize}
    \item To what extent were the thesis objectives achieved?
    \item Robot model validated against Klein (2023); deviations explained by simplified tendon model
    \item RL pipeline functional for both reach and place; tendon-driven variant achievable but with different convergence characteristics
\end{itemize}

\section{Tendon Wrist Transferability}
\begin{itemize}
    \item Interpretation of wrist transfer results
    \item Under what conditions does the tendon wrist improve or degrade performance?
    \item Implications for robot design: when is a compliant wrist advantageous?
\end{itemize}

\section{Limitations}
\begin{itemize}
    \item Simplified tendon model (PD drives, no cable stretch or backlash)
    \item Rigid objects only; no deformable cloth tested
    \item Simulation-only results; no sim-to-real validation
    \item Wrist experiment limited to one standard arm type
\end{itemize}

\section{Practical Implications}
\begin{itemize}
    \item Modular simulation pipeline reusable for future tasks
    \item Workspace analysis informs physical task design
    \item Wrist results guide hardware decisions for the master thesis
\end{itemize}

%--- ~3 pages ---%
\chapter{Conclusion}\label{ch:conclusion}
\markboth{Conclusion}{}

\section{Conclusion}
\begin{itemize}
    \item Recap accomplishments: validated simulation, RL-controlled tendon robot, wrist transferability assessment
    \item Answer the research question concisely
    \item Highlight key quantitative results
\end{itemize}

\section{Outlook}
\begin{itemize}
    \item Direct next steps (master thesis scope):
    \begin{enumerate}
        \item Conveyor-based sorting with selectivity (T-shirt vs.\ non-target)
        \item Deformable cloth simulation integration
        \item Second robot arm for stretching and inspection
        \item Camera-based damage classification
        \item Multi-agent coordination
        \item ROS 2 integration for sim-to-real transfer
    \end{enumerate}
    \item More realistic tendon model (cable dynamics, backlash)
    \item Domain randomization for transfer robustness
\end{itemize}
